\chapter{Evaluation Methodology and Research Rigour}\label{ch:evaluation}

Your thesis already showed good instincts: stratified cross-validation, a test
set touched exactly once, and ties broken on generalisation. This chapter
carries the same discipline into the hardware project. It covers what to
measure, what to compare against, how to tell a real difference from noise, and
how to make every number reproducible.

% =========================================================================
\section{What to measure}\label{sec:eval:metrics}

\begin{table}[htbp]
\centering\footnotesize
\caption{Metric catalogue. Use the same definitions everywhere in the project.}
\label{tab:eval:metrics}
\begin{tabularx}{\textwidth}{@{}lLL@{}}
\toprule
Group & Metric & Notes \\
\midrule
Speech quality & PESQ (wideband)~\cite{itu2007pesqwb}, STOI~\cite{taal2011stoi}, SI-SDR~\cite{leroux2019sisdr}, $\Delta$SNR & Against the clean reference, after alignment; per utterance; PESQ at \qty{16}{\kHz} \\
Implementation fidelity & Bit-exactness against golden; output SNR of fixed point against float & Separates \enquote{the hardware is wrong} from \enquote{the precision is too low} \\
Intent & Macro-F1, per-class F1, severity MAE & Thesis split; test set used once \\
Timing & Cycles per frame (max and mean), algorithmic and end-to-end latency (ms) & Measured, not estimated, by M3 \\
FPGA cost & LUT, FF, DSP, BRAM; WNS at the target clock & From Vivado reports \\
ASIC cost & Cell and core area (mm$^2$), setup and hold slack & From LibreLane reports \\
Power and energy & Dynamic and static power (mW); energy per frame (\textmu J) & State the activity source \\
Memory & State and table size (KiB); accesses per frame & From the spec and RTL \\
\bottomrule
\end{tabularx}
\end{table}

\begin{conceptbox}{Know what your metrics can and cannot say}
PESQ and STOI were designed and validated for particular conditions. They
predict average listener judgements for normal-hearing listeners on telephone
and enhancement tasks. Neither is a perfect proxy for what a hearing-aid user
perceives. Use them to compare \emph{configurations of the same algorithm}, which
is exactly what RQ1 and RQ2 need, rather than to claim absolute perceptual
quality.
\end{conceptbox}

% =========================================================================
\section{Baselines}\label{sec:eval:baselines}

Every result needs a reference point:
\begin{itemize}
  \item \textbf{Unprocessed noisy input}: does enhancement help at all?
  \item \textbf{Float64 thesis model}: the reference for RQ1 and RQ2.
  \item \textbf{Float32 thesis model}: a sanity check of how much precision
        actually matters.
  \item \textbf{Standard Wiener filter} ($\beta=1$, no tilt): isolates the
        effect of the parametric control.
  \item For intent: \textbf{the thesis's mpnet + SVM} at macro-F1 0.957 and MAE
        0.100.
\end{itemize}

% =========================================================================
\section{Experimental design}\label{sec:eval:design}

\begin{itemize}
  \item \textbf{Development versus test.} Make every choice (word lengths, table
        sizes, frame configuration, quantisation scales) on the development
        subset. Run the test set once, at the end, for the reported numbers.
        This is the same rule as your thesis's test-once discipline.
  \item \textbf{One factor at a time, then jointly.} Single-factor sweeps find
        each knee. A joint run checks that the errors do not add up beyond the
        target.
  \item \textbf{Report by condition.} VoiceBank+DEMAND's test set has four input
        SNRs. Effects often differ between low and high SNR, so report each
        condition as well as the overall mean.
  \item \textbf{One change per experiment.} Each experiment folder records
        exactly one hypothesis, one configuration and one commit.
\end{itemize}

% =========================================================================
\section{Statistics: is the difference real?}\label{sec:eval:stats}

Most comparisons in this project are \emph{paired}: the same 824 utterances,
processed two ways. Use that pairing:
\begin{enumerate}
  \item compute the per-utterance difference $d_i$, for example
        $\text{PESQ}_\text{fixed}-\text{PESQ}_\text{float}$;
  \item report the mean $\bar d$ with a 95\% \emph{bootstrap} confidence
        interval: resample the 824 differences with replacement many times and
        take the 2.5th and 97.5th percentiles of the means;
  \item for \enquote{is there any difference at all?}, use a Wilcoxon
        signed-rank test on the $d_i$;
  \item for RQ1's equivalence question, check that the \emph{whole} confidence
        interval lies within the tolerance ($|\Delta\text{PESQ}|\le0.1$). This is
        stronger, and more honest, than \enquote{not significantly different}.
\end{enumerate}

\begin{lstlisting}[style=py,caption={Paired bootstrap confidence interval for a mean difference.}]
import numpy as np

def paired_ci(a, b, n_boot=10_000, seed=0):
    d = np.asarray(a) - np.asarray(b)
    rng = np.random.default_rng(seed)
    idx = rng.integers(0, len(d), size=(n_boot, len(d)))
    means = d[idx].mean(axis=1)
    return d.mean(), np.percentile(means, [2.5, 97.5])
\end{lstlisting}

For the NLP results, the test set has only about 166 sentences. Report how many
predictions changed between float and int8, not just the change in F1. A
difference of one or two sentences is within noise.

% =========================================================================
\section{Reproducibility}\label{sec:eval:repro}

Follow the \enquote{good enough practices} of scientific
computing~\cite{ch12-13:wilson2017goodenough,ch12-13:sandve2013tenrules}:
\begin{itemize}
  \item Every experiment has a configuration file, a fixed seed and a commit
        hash, stored with its raw results (CSV), not just the plots.
  \item One command per figure. A \code{make figures} target, or a script,
        regenerates every table and plot from raw results.
  \item Record tool versions: Python packages (lock file), Verilator, cocotb,
        Vivado, the PYNQ image, LibreLane and the PDK version.
  \item Tag the repository at every milestone.
  \item Raw audio stays out of git, but record how to obtain it (the dataset DOI
        and a script).
\end{itemize}

% =========================================================================
\section{Presenting results}\label{sec:eval:present}

Plan your key tables and figures now, and fill them as you go:
\begin{itemize}
  \item quality against word length, per signal and joint (RQ1);
  \item quality against latency, with state memory (RQ2): the Pareto plot;
  \item resources and area against word length, next to the quality plot (RQ3);
  \item energy per frame for the architectures you compared (RQ3, H4);
  \item macro-F1 against model size and precision (RQ4).
\end{itemize}

\begin{table}[htbp]
\centering\footnotesize
\caption{Template for the headline results table (fill in at M5).}
\label{tab:eval:headline}
\begin{tabularx}{\textwidth}{@{}LCCCCC@{}}
\toprule
Configuration & PESQ & STOI & Latency & State & Energy/frame \\
\midrule
Noisy input & & & --- & --- & --- \\
Float64 thesis (48\,kHz, 2048 FFT) & & & 21.33\,ms & 23,575 words & --- \\
Fixed point, chosen spec (golden = RTL) & & & & & \\
FPGA (measured) & = above & = above & & & \\
ASIC (Track B, estimated) & = above & = above & & & \\
\bottomrule
\end{tabularx}
\end{table}

% =========================================================================
\section{Threats to validity}\label{sec:eval:threats}

A strong report names its own weaknesses before a reviewer does. Discuss at
least these:
\begin{itemize}
  \item \textbf{Objective metrics are not perception}, least of all for
        hearing-impaired listeners. The thesis's open calibration problem
        remains open.
  \item \textbf{The data is not hearing-aid data.} VoiceBank+DEMAND is studio
        speech with added noise, not a hearing-aid microphone in a real room.
  \item \textbf{The intent data is typed and LLM-generated}, as the thesis
        noted, and the TTS data in Route C is not real speech.
  \item \textbf{FPGA power is not ASIC power}, and \textbf{130\,nm open-PDK
        results are not modern hearing-aid silicon}.
  \item \textbf{Estimates are not measurements.} Every number has a label.
\end{itemize}

% =========================================================================
\section{Ethics and data}\label{sec:eval:ethics}

If you record anyone's voice, for latency tests or Route C, get informed
consent, say what the recordings are for, store them securely, and delete them
when the project ends. Recordings of identifiable people are personal data under
GDPR. Check your university's guidance before recording anyone. Keep public
datasets under their licences, and cite them.

\begin{donebox}[Your evaluation is sound when]
\begin{checklist}
  \item Every number in the report has a definition, a method, a dataset split and a commit.
  \item Every comparison has a baseline and a confidence interval.
  \item The test set was used once, at the end.
  \item Every figure regenerates from raw results with one command.
  \item The threats to validity are written down.
\end{checklist}
\end{donebox}
